\section*{Chapter \ref{ch:ll:queues-ring-buffers-and-shared-memory} --- Queues, Ring Buffers and Shared Memory}

\begin{solution}{exo:ll:queues-ring-buffers-and-shared-memory:1}
$70\,001 - 69\,500 = 501$ messages wait. The next write goes to slot $70\,001 \bmod 1\,024 = 369$. The ring is full at 1\,024 waiting,
so $1\,024 - 501 = 523$ more.
\end{solution}

\begin{solution}{exo:ll:queues-ring-buffers-and-shared-memory:2}
$2^{64}/(2.5 \times 10^8)$ seconds is about 2\,338 years: never. A 32-bit position wraps in 17.2 seconds, which is why positions are
64 bits and only the slot index is taken modulo the size.
\end{solution}

\begin{solution}{exo:ll:queues-ring-buffers-and-shared-memory:3}
$0.8^n < 10^{-6}$ for $n \ge 62$; the next power of two is 64 slots.
\end{solution}

\begin{solution}{exo:ll:queues-ring-buffers-and-shared-memory:4}
Reading the other side's position means reading a line the other core writes at every message, so it moves between the cores each
time. With the copies, each side reads the other's line only when the ring looks full or empty; without them every message adds one
or two line transfers, tens of nanoseconds under contention.
\end{solution}

\begin{solution}{exo:ll:queues-ring-buffers-and-shared-memory:5}
The backlog grows at $3\,000\,000 - 400\,000 = 2.6$ million a second and reaches 131\,072 after 50.4 milliseconds: the monitor is
lapped about 30 milliseconds before the burst ends.
\end{solution}

\begin{solution}{exo:ll:queues-ring-buffers-and-shared-memory:6}
In the C++ memory model a non-atomic read that races with a write is a \omterm{def:ll:rust-for-low-latency:race}{data race}, which is \omterm{def:ll:cpp-for-latency-iii-the-compiler:ub}{undefined behaviour} whatever is done with
the value afterwards; the compiler may assume it does not happen. Relaxed atomic loads and stores of the words make the race defined,
and the sequence check makes the result consistent.
\end{solution}

\begin{solution}{exo:ll:queues-ring-buffers-and-shared-memory:7}
Both positions then share a line that the producer writes at every message and the consumer at every message: \omterm{def:ll:memory-hierarchy-and-caches:false}{false sharing}
(chapter~3). Throughput should fall, since every write invalidates the other core's copy of the line, and the round trip changes
less, since in a ping-pong the lines move anyway; the measurement decides by how much.
\end{solution}

\begin{solution}{exo:ll:queues-ring-buffers-and-shared-memory:8}
A pointer is an address in one process's address space and means nothing, or something else, in the other; a
\texttt{std::string} holds such a pointer to memory that is not in the region at all. Shared layouts hold offsets and fixed-size
fields only (a fixed character array for the name), with a version number.
\end{solution}

\begin{solution}{pb:ll:queues-ring-buffers-and-shared-memory:1}
\begin{enumerate}
\item $65\,536 \times 64 = 4$~MiB plus the header: twice a 2~MiB second-level cache; the ring's hot part is its recent slots.
\item 7.5\%, 12\% and 20\%.
\item The monitor, the only consumer slower than 2 million a second.
\item The exchange will not wait: a producer that blocked on its slowest reader would drop packets at the socket instead.
\item $65\,536/(2\,000\,000 - 1\,500\,000) = 131$ milliseconds.
\item Yes, after 131 milliseconds of a 500-millisecond burst.
\item At least $0.5 \times 500\,000 = 250\,000$ slots: 262\,144 as a power of two.
\item 16~MiB.
\item 5\,000 lines of 64 bytes, 320\,000 bytes; at 2 million updates a second against 1.5 million reads, it skips a quarter of them,
seeing only the latest.
\item It reports the gap, jumps to a recent position and rebuilds its view (from its own snapshot or the book builder's).
\item Losing market data is recoverable (snapshots, \omterm{def:ll:queues-ring-buffers-and-shared-memory:bp}{conflation}) and waiting is not; losing or reordering an order is neither, so the
order path waits or refuses.
\item $\P(N \ge n) = 0.12^n$, negligible for any reasonable $n$; optimistic because arrivals come in bursts far from Poisson.
\item \textbf{Named result.} The monitor is lapped 131 milliseconds into a 2-million-a-second burst on a 65\,536-slot ring; reading
through per-instrument \omterm{def:ll:queues-ring-buffers-and-shared-memory:seqlock}{sequence locks} instead, it skips 25\% of the updates and is never lapped.
\item About \qty{200}{\nano\second} per round trip on this laptop: a consumer adds a cursor, not a queue, and its cost is its own reading.
\item Because two binaries built at different times, or in different languages, read it: a change must be detected, not misread.
\item Hundreds of times the ring's latency per hop, with a tail in milliseconds, as the published comparison found.
\item Both on the network card's node, on cores that hand lines over fastest (chapter~4).
\item Record the arrival rate per millisecond from the capture at the wire (chapter~23) and keep the distribution of burst lengths
and peaks.
\item Invalidate the affected state and wait for, or request, a resynchronisation, as the book builder does on a gap (chapter~19).
\item The bursts it must absorb.
\end{enumerate}
\end{solution}

\begin{solution}{iq:ll:queues-ring-buffers-and-shared-memory:1}
A power-of-two array, a producer position and a consumer position on separate \omterm{def:ll:memory-hierarchy-and-caches:line}{cache lines}, each written by one side. The producer
writes the slot, then stores its position with release; the consumer loads it with acquire, reads, and stores its own with release,
which the producer loads with acquire before reusing slots. Cache the other side's position.
\iqlookfor{release/acquire pairs, separate lines, cached positions.}
\end{solution}

\begin{solution}{iq:ll:queues-ring-buffers-and-shared-memory:2}
No lock to contend for and no sleeping: the consumer spins on a position in its cache and sees the producer's write as a \omterm{def:ll:memory-hierarchy-and-caches:line}{cache
line} transfer, a hundred nanoseconds or so, where a condition variable puts the thread to sleep in the kernel and wakes it with a
system call and a context switch, tens of microseconds.
\iqlookfor{kernel involvement versus cache-coherence cost.}
\end{solution}

\begin{solution}{iq:ll:queues-ring-buffers-and-shared-memory:3}
A counter made odd during a write and even after; readers copy and retry if the counter changed or was odd. For small values with
one writer and many readers who want only the latest (top of book). Pitfalls: the copy must use atomic loads in C++ to avoid a \omterm{def:ll:rust-for-low-latency:race}{data
race}; readers can starve under constant writes; the writer must be single.
\iqlookfor{the retry protocol and the language-level \omterm{def:ll:rust-for-low-latency:race}{data race}.}
\end{solution}

\begin{solution}{iq:ll:queues-ring-buffers-and-shared-memory:4}
One broadcast ring written by the feed handler, each consumer at its own position, the producer never waiting; consumers that
cannot keep up get conflated views (latest value per instrument) or detect overruns and resynchronise from a snapshot; sizes set
from recorded bursts.
\iqlookfor{a policy for the slowest consumer.}
\end{solution}

\begin{solution}{iq:ll:queues-ring-buffers-and-shared-memory:5}
Pointers and anything that contains them (strings, containers, virtual tables), fields whose size or alignment depends on the
compiler, and layouts without a version; locks that cannot be recovered if the holder dies.
\iqlookfor{address-space independence and versioning.}
\end{solution}

\begin{solution}{iq:ll:queues-ring-buffers-and-shared-memory:6}
Start from the M/M/1 tail $\P(N \ge n) = \rho^n$ for a target overflow probability; then check against recorded bursts, which are far
from Poisson and last seconds; size for the worst burst plus margin, and design what happens on overflow. Queueing theory gives the
order of magnitude and the effect of utilisation; it misses correlation and bursts.
\iqlookfor{theory for the shape, data for the size.}
\end{solution}
